\subsection{Additional controls and sensitivity analyses}
\label{app:revision-controls}

We varied the persistence of the intervention, the boundary response,
and the observation window to identify which parts of the model
supported the transition. These controls were added in response to
review and retained the original initialization and 1,000-step baseline.
Unless specified otherwise, they used $L=64$, $\lambda=0.02$,
$h=\infty$, and 3,000 post-intervention steps. Each new rule or pulse
arm had a matched arm that retained the baseline rule. We required
exact pre-intervention agreement of their lattice and Sell histories,
and checked the price prehistory for financial arms. Held-boundary
experiments needing only lattice observables used an allocation-free
kernel with the production setup draws and rule tables. Tests matched
its full Sell history and final grid to the market simulator, and
1,800 overlapping archived seed/fraction outcomes also agreed. The
reduced kernel reproduced the local dynamics when prices could not
feed back, but did not calculate financial observables.

\paragraph{Matched transient perturbation and rule screen.}
The matched pulse placed targeted traders in Sell once while leaving
their subsequent response rule unchanged
(Supplementary Table~\ref{tab:revision-pulse}). At activation, we wrote
Sell at the rule-switch target sites before orders executed and retained
Rule~19280 everywhere. The pulse could therefore change the first quote,
whereas persistent activation preserved that quote and changed later
local updates. Fractions matched the number of targeted sites; sites
already in Sell did not change state. For $f=0.15,0.20,0.25$ under both
boundaries, no pulse reached a majority in 50 seeds (0/50; Wilson
upper bound 0.071). Persistent Rule~19286 activation yielded 0/50,
50/50, and 50/50, respectively. The Wilson lower bound for 50/50
was 0.929. A one-time increase in Sell occupancy at these targets
was thus insufficient to reproduce the collective response to a
lasting rule change. The comparison applied to this matched pulse;
other choices of state perturbation could have different outcomes.

The rule screen tested all 18 ways to change one entry of Rule~19280
(Supplementary Table~\ref{tab:revision-rules}). Each of the nine
neighbor sums had two alternative outputs. We activated each resulting
rule at the same three fractions using 20 seeds under both boundaries.
At $f=0.20$, six rules produced majority crossings in all 20 seeds
under each boundary, while the other twelve produced none. Effective
changes included entries $c=2,3,4,5$ as well as the selected $c=1$
change. Every change that replaced a Buy or Hold response with Sell at
a neighbor sum from one to five produced a majority, as did the change
from Buy to Hold at a sum of three; the same Buy-to-Hold change at a
sum of one or two, every change that moved a response away from Sell,
and every change at a sum of six or more produced none. The input being altered and its replacement output therefore
mattered even when the number of changed entries was fixed. The screen
was descriptive and used no adjustment for multiple comparisons.
An observed 0/20 had a Wilson upper bound of 0.161, and 20/20 had
a lower bound of 0.839. The screen identified other rules capable
of producing majorities over the tested window without defining
a universal class of activation thresholds.

\paragraph{Frozen band and non-monotone dealer response.}
Boundary decision widths changed as the quote moved, allowing the
same price departure to elicit different trader responses over time
(Supplementary Table~\ref{tab:revision-band}). We removed this adaptation
by freezing each path's decision deviation immediately before activation.
Traders retained their $k_j$ values and the moving center continued to
update. The half-width $k_j\sigma_{1000}$ then stayed constant within
a path, although the diagnostic deviation still evolved. Each frozen
arm shared its adaptive arm's exact prehistory, so widths could differ
across seeds and capacities. At $f=0.175$, the $h=0.1$ minus
$h=\infty$ majority contrast changed from 0.38 with adaptation
to -0.02 with frozen widths. The paired difference of these
contrasts was 0.40, with a 95\% seed-bootstrap interval
of [0.16, 0.64]. Removing width adaptation therefore removed
the low-capacity advantage in the point estimates. Boundary Sell
occupancy also changed, but the frozen capacity response remained
non-monotone and individual probability intervals were wide. Width
adaptation contributed to the observed advantage without isolating
a unique microscopic cause of the entire capacity response.

\paragraph{Price-impact sensitivity.}
The price-impact coefficient set the quote change produced by a given
executed imbalance (Supplementary Table~\ref{tab:revision-impact}).
At $f=0.20$, we varied $\lambda$ over $\{0.005,0.02,0.05\}$ and
capacity over $h\in\{0.1,1,\infty\}$. A smaller coefficient reduced
the immediate price movement from a given imbalance. Under the endogenous
boundary, the changed quote also entered the decision bands that set
perimeter states. All 50 seeds transitioned in every cell. Median
drawdown increased with capacity at each coefficient, although its
magnitude changed by a factor of seven to ten across the coefficient
range. Dealer resources thus
limited the price decline across a decade of impact choices, while
collective selling persisted at this activation fraction throughout
the sweep. These uncalibrated coefficients describe model sensitivity;
they do not convert the simulated drawdowns into estimates for an
empirical market. Both the amount of eligible flow the dealer could
absorb and the impact assigned to filled imbalance determined the
financial magnitude of the transition.

\paragraph{Horizon scaling and transition criterion.}
A longer observation window gave the altered local dynamics more time
to reach the selected Sell fraction
(Supplementary Table~\ref{tab:revision-horizons}). Held-boundary checks
used 50 seeds on $f=0.140,0.145,\ldots,0.250$, extended to 0.300 at
$L=32$ to bracket the short-horizon response. We compared the original
3,000 steps with $T_{\rm post}=3000(L/64)^z$ for $z=1,2$ at the
endpoint sizes. These windows were 1,500 and 750 steps at $L=32$,
and 6,000 and 12,000 at $L=128$. Both scalings preserved the ordering
of endpoint midpoints. At $L=64$, increasing the horizon from 3,000
to 9,000 steps lowered the majority midpoint from about 0.178 to 0.167.
Longer endpoint runs likewise continued to move the curve. Changing
the crossing criterion to 0.4 or 0.6 retained the endpoint ordering
and downward drift with time. At $L=64$ and 3,000 steps, the three
criteria gave midpoints of approximately 0.172, 0.178, and 0.186.
The qualitative size response persisted, but the numerical midpoint
depended on both the criterion and horizon. Neither these comparisons
nor the final 500-step mean established equilibration at any size.

\paragraph{Baseline characterization and long no-switch paths.}
The pre-switch paths remained mixed, with similar mean Sell occupancy
in the two windows preceding activation. Across 50 endogenous
unlimited-dealer seeds, the mean interior Sell fraction was
0.2494 over steps 501--750 and
0.2503 over steps 751--1,000. In the
latter window, the median within-seed return standard deviation was
$0.0003$, and the median within-seed
1st and 99th percentiles were $-0.00068$
and $0.00066$. The mean zero-return
fraction was 0.006. These summaries
characterized the starting ensemble without establishing stationarity
or empirical calibration. We extended the first ten production seeds
under each boundary for 100,000 additional no-switch steps. None
crossed a majority (0/10 per boundary; Wilson upper bound 0.278
for the ensemble event probability). The extension followed existing
seeds and was not an independent replication. It showed that those
paths remained below a majority over that longer interval, while
leaving global stability and the distinction between long transients
and attractors unresolved.

\paragraph{Fit and grid uncertainty.}
A monotone logistic curve supplied a second estimate of transition
midpoint and width alongside empirical interpolation
(Supplementary Table~\ref{tab:revision-fits}). Interpolation used
the first upward crossing of each target probability and allowed
local decreases in the sampled curve. The logistic fit imposed a
monotone response as a working shape model. Both analyses used 2,000
bootstrap replicates that resampled whole seeds across fractions.
Keeping each seed together preserved the dependence arising from
shared initialization, so standard errors reflected seed variation
without treating fractions as independent observations. Each
susceptibility replicate reselected the maximum over the sampled grid.
Its location remained a sampled fraction with finite local resolution.
The resulting intervals measured uncertainty conditional on the refined
grid and observation horizon. They did not account for discretization
error, adaptive-grid selection uncertainty, or changes caused by
extending the horizon. No susceptibility exponent was estimated from
these finite-size summaries.
